\documentclass[9pt]{extarticle}
\usepackage[T1]{fontenc}
\usepackage[utf8]{inputenc}
\usepackage{lmodern}
\usepackage[top=0.8in,bottom=0.8in,left=0.65in,right=0.65in]{geometry}
\usepackage{microtype}
\usepackage{amsmath,amssymb}
\usepackage{graphicx}
\usepackage{booktabs,longtable,array,calc,etoolbox}
\usepackage{url}
\usepackage{multicol,caption}
\usepackage[hidelinks]{hyperref}
\DeclareUnicodeCharacter{2212}{\ensuremath{-}}
\DeclareUnicodeCharacter{2265}{\ensuremath{\geq}}
\DeclareUnicodeCharacter{2192}{\ensuremath{\rightarrow}}
\DeclareUnicodeCharacter{03BA}{\ensuremath{\kappa}}
\DeclareUnicodeCharacter{03B4}{\ensuremath{\delta}}
\DeclareUnicodeCharacter{0302}{\^{}}
\providecommand{\tightlist}{\setlength{\itemsep}{0pt}\setlength{\parskip}{0pt}}
\setcounter{secnumdepth}{0}  % section numbers are part of the heading text
\setlength{\parskip}{0.3em}\setlength{\parindent}{0pt}
\emergencystretch=2em
\setlength{\columnsep}{0.25in}
\captionsetup{font=small,skip=2pt}
\AtBeginEnvironment{longtable}{\scriptsize\raggedright}
% the manuscript supplies its own "References" heading above the bibliography
\patchcmd{\thebibliography}{\section*{\refname}}{}{}{\PackageWarning{main}{bibliography heading not patched}}
\title{Supplement to: What Reached the Executor? Six Candidate Measurement Checks for Evaluating Runtime Defenses of LLM Agents: A Case Study with a Negative Result}
% AUTHORS: names as given by the owner. Spelling of the second name is unconfirmed (given as 'RezaManzour').
% Emails as given by the owner. Affiliation (independent) as given by the owner. Corresponding author (Matin Shirazi, the first author) is a default the authors may change: fill them in before posting.
\author{Matin Shirazi (Shirazimatin@gmail.com) \and Reza Manzour (Rezamanzourolajdad@gmail.com)\\[0.4em]
{Independent Researchers}\\
{Corresponding author: Matin Shirazi}}
\date{}
\begin{document}
\maketitle
\noindent\textit{Supplement to the main paper. Section numbers (\S) refer to the main paper. The text below is reproduced unchanged from the same manuscript source.}

\section{S1 Subsections moved from the main paper}

\subsection{5.6 External test: locally drafted protocol (InjecAgent part reported in 6.6)}

Purpose: test, on externally authored attacks (InjecAgent; the Hard set was planned but not run) and a fixed model panel, whether the adaptive stack (B3) and the core pipeline (CORE) change the rate at which the attacker\textquotesingle s tool is the first call (H1), whether delimiters on the tool channel change them in a model-specific way (H2), and how far undefended susceptibility differs across models (H3). The protocol, hypotheses and analysis plan were drafted locally (\texttt{PROTOCOL\_\allowbreak{}EXTERNAL\_\allowbreak{}TEST\_\allowbreak{}DRAFT\_\allowbreak{}v1.md}, labeled not frozen) and were not externally registered; §6.6 gives the chronology that the repository supports.

\emph{Panel.} Four open targets (llama-4-maverick, deepseek-v4.1-flash, qwen3.8-flash, glm-4.7), one closed target (gpt-5.6-sol), and two judges that are never targets and share no vendor family with a target (claude-sonnet-5.5, grok-4.7), used only on the judge-scored tracks. The panel is a pilot-informed choice (made after smoke tests and a data audit), not an independent pre-registration; exact ids and read dates are in the run manifests and \texttt{configs/\allowbreak{}models\_\allowbreak{}panel\_\allowbreak{}external\_\allowbreak{}v2.yaml}.

\emph{Data.} (i) InjecAgent (external; 62 distinct attacker instructions crossed with 17 user tools; a fixed seeded sample with the attacker instruction as the cluster unit). (ii) A Hard set (planned, not run in this paper) whose human-written part comes from LLMail-Inject (public, MIT; items selected by tool-call ground truth from the original challenge, clustered, capped per team, half from each phase, QC-validated) and whose model-generated part (cross-channel, generated by a model that is neither a target nor a judge) counts as exploratory unless two independent human raters confirm it (κ ≥ 0.70, agreement ≥ 95\%). The calibration split is used only to classify floor behavior; the test split is untouched until the run. Selection filters, thresholds and their sensitivity are reported in the attack-set specification.

\emph{Arms and endpoint.} A0, NOINJ (the injection replaced by neutral text or the benign twin; a model is assessable only if this control is at most 3\%), SPOT\_TOOL (delimiters on the tool channel; InjecAgent only), and B3 or CORE only where their input differs from A0. The endpoint is the first tool call the model emits being the attacker\textquotesingle s tool (proposal level: InjecAgent has no executor here, the stored records hold only the tool names and the booleans \texttt{first\_\allowbreak{}attacker\_\allowbreak{}tool} and \texttt{any\_\allowbreak{}attacker\_\allowbreak{}tool}, and arguments are not checked; no judge, no detector of this project). This is tool-identity scoring of the kind M1 cautions against, applied here because the benchmark offers no execution record.

\emph{Analysis.} Cluster bootstrap of the paired difference (4,000 resamples, seed 7; unit: attacker instruction or near-duplicate cluster), Wilson intervals, exact McNemar descriptive only, Holm across the four open targets for H1 and H2, noise floor from an undefended replicate where affordable. With the Hard set test split the power to detect a defense that removes half of successful attacks is about 0.6 at an undefended rate of 0.2 and above 0.95 at 0.4 (alpha = 0.0125), so models at the floor are described, not tested.

\emph{Scope under budget.} The run is sized to the project\textquotesingle s compute budget: InjecAgent (A0 and SPOT\_TOOL on 124 cases, NOINJ on 40) on the four open targets and A0 on 40 cases for the closed target; Hard set test (planned, not run: A0 on the full human-written test split plus the generated test items, NOINJ on 40) on the four open targets and A0 on 40 items for the closed target. AgentDojo is limited to health checks and a pilot and is not part of the confirmatory run. Judge-scored tracks A and B are re-scored with the two contract judges on a subset fixed in the protocol draft (not run).

\subsection{6.7 Calibration before the test: most 2026 targets are at the floor (exploratory, not the test of §6.6)}

Before spending the budget on the InjecAgent run of §6.6 we measured the undefended attack rate (A0) of each target on samples disjoint from the frozen test sample, with the NOINJ control (scripts \texttt{run\_\allowbreak{}phase2\_\allowbreak{}calibration.py} and \texttt{run\_\allowbreak{}injecagent\_\allowbreak{}panel.py}; records in \texttt{experiments/\allowbreak{}external/\allowbreak{}}). Exact 95\% upper bounds for zero events are 8.8\% at n = 40 and 5.3\% at n = 68.

\begin{longtable}[]{@{}
  >{\raggedright\arraybackslash}p{(\columnwidth - 8\tabcolsep) * \real{0.2000}}
  >{\raggedright\arraybackslash}p{(\columnwidth - 8\tabcolsep) * \real{0.2000}}
  >{\raggedright\arraybackslash}p{(\columnwidth - 8\tabcolsep) * \real{0.2000}}
  >{\raggedright\arraybackslash}p{(\columnwidth - 8\tabcolsep) * \real{0.2000}}
  >{\raggedright\arraybackslash}p{(\columnwidth - 8\tabcolsep) * \real{0.2000}}@{}}
\toprule\noalign{}
\begin{minipage}[b]{\linewidth}\raggedright
target
\end{minipage} & \begin{minipage}[b]{\linewidth}\raggedright
InjecAgent A0 (n = 40)
\end{minipage} & \begin{minipage}[b]{\linewidth}\raggedright
Hard set, human-written A0 (n = 68)
\end{minipage} & \begin{minipage}[b]{\linewidth}\raggedright
NOINJ
\end{minipage} & \begin{minipage}[b]{\linewidth}\raggedright
classification
\end{minipage} \\
\midrule\noalign{}
\endhead
\bottomrule\noalign{}
\endlastfoot
llama-4-maverick & 5/40 & 1/68 & 0 & measurable on InjecAgent; floor on Hard set \\
qwen3.8-flash & 6/36 scored (4 provider errors) & 0/40 scored (28 provider errors) & 0 & measurable on InjecAgent (limited n); Hard set undetermined \\
glm-4.7 & 1/40 & 1/68 & 0 & floor \\
deepseek-v4.1-flash & 0/40 (upper bound 8.8\%) & 0/68 (upper bound 5.3\%) & 0 & floor \\
gpt-5.6-sol (closed) & 0/40 (upper bound 8.8\%) & 0/20 & 0/40 & floor on InjecAgent; Hard set undetermined (n \textless{} 40) \\
\end{longtable}


Two readings follow. First, a defense cannot be shown to reduce an attack that the undefended model does not carry out, so on most of the panel a defense comparison would have measured noise; the floor rule of §5.6 therefore decides which targets enter the comparison of §6.6 (llama-4-maverick and qwen3.8-flash on InjecAgent). Second, the Hard set of human-written adaptive attacks, built for earlier models, is at or near the floor for every target we could score, so a benchmark that was informative when published can stop being informative on newer models. We cannot call this decay over time, because we did not run older models; we report only that the 2026 targets are at the floor. Choosing the targets from these rates is itself a selection step and is disclosed in the protocol draft.

\subsection{8.5 Artifact and process limitations}

Test suite status: an internal project review on 2026-09-30 recorded 283 passing tests on \texttt{main} and 23 failures on the working branch that contains the harness (note in \texttt{docs/\allowbreak{}research/\allowbreak{}}); the full suite on the working branch on 2026-10-01 gave 661 passed, 32 failed and 6 skipped (\texttt{docs/\allowbreak{}research/\allowbreak{}TEST\_\allowbreak{}SUITE\_\allowbreak{}STATUS\_\allowbreak{}20261001.md}). The two figures come from different branches and dates and are not directly comparable. On the public tree of this revision (2026-10-07), with \texttt{src/\allowbreak{}} on the Python path, the full suite gave 283 passed and 3 failed with \texttt{openai} and \texttt{google-genai} not installed, and 287 passed and 0 failed once \texttt{httpx}, \texttt{scikit-learn}, \texttt{openai} and \texttt{google-genai} were installed; without \texttt{src/\allowbreak{}} on the path two further tests fail because \texttt{adapti\_\allowbreak{}guard} is not installed as a package (\texttt{docs/\allowbreak{}research/\allowbreak{}REPRODUCIBILITY\_\allowbreak{}REPORT\_\allowbreak{}20261007.md}). In scratch checkouts of the two archive commits that contain the harness (\texttt{1ae0fb4d}, \texttt{dfbea801}) the suite gave 670 passed and 24 failed (the same 24 in both), 15 of them campaign-preflight tests that read a \texttt{PENDING\_\allowbreak{}BUDGET\_\allowbreak{}APPROVAL} placeholder as an integer; the earlier figures of 661 passed and 32 failed were not reproduced. The failures are listed in the reproducibility report and were not repaired. The recorded causes include old campaign-test configuration, missing optional packages and some undiagnosed assertion failures; we did not repair them. The tests that bind the manuscript (number ledger, panel registry, external adapters, attack sets) pass. A statement in an earlier internal report that the Phase-1 detector was authored before the holdout pack cites commit \texttt{c462945}; that commit is on the public remote (reachable from pull-request refs) and its rewritten equivalent \texttt{12414c8} is in the history of \texttt{main}, and the recorded commit order supports the statement (§4.4). Amendment 10 and the fixed-tool-policy arms are proposals awaiting the authors\textquotesingle{} approval. Live spend by experiment: E2 and E3 about \$0.23 (runner-reported, per-run \texttt{cost\_\allowbreak{}summary.json}), the E4 channel re-run \$0.0669, and the InjecAgent run of §6.6 about \$0.21 (sum of per-episode cost records); calibration (§6.7) and pilot (Appendix A) spend is not included in these figures. Results depend on this small scale and should be repeated at larger K before being relied on. AI assistance in producing repository artifacts is disclosed in §10.

\textbf{Provenance of the E2 and E3 runs.} For the four multi-turn harness runs of 2026-09-30 (\texttt{HARNESS\_\allowbreak{}V2\_\allowbreak{}\{EXPLORATORY\_\allowbreak{}SMOKE,\allowbreak{}EXPLORATORY,\allowbreak{}INDEPENDENT\_\allowbreak{}SCREEN,\allowbreak{}INDEPENDENT\_\allowbreak{}DEFENDED\}\_20260930}), no branch or tag of the repository contains pre- or post-run provider-key usage snapshots (\texttt{/\allowbreak{}auth/\allowbreak{}key}) or a \texttt{progress.log}. The runner and repository commits recorded in their manifests are not reachable from any branch or tag of the repository. They could still be fetched from the public remote by full SHA (checked on 2026-10-03), but the archive tag \texttt{case-study-v1} named in \texttt{REPRODUCIBILITY.md} does not exist on the remote, the original-history bundle is held by the authors, and unreachable commits are not guaranteed to remain retrievable. No per-request cost records are in the tree of any branch or tag. The per-request logs exist only in the unreachable archive commit and match the runner\textquotesingle s \texttt{cost\_\allowbreak{}summary.json} totals for three runs (\$0.0941 against \$0.0944 for the defended run); none was compared with provider-side usage, so the reported spend rests on the runner\textquotesingle s own records. No record of the authors\textquotesingle{} approval of these runs is preserved in the repository, its pull requests or its commit messages.

\textbf{Non-delivered episodes and provider errors in the E3 pairing.} The paired E3 comparisons (§6.3, Appendix D) pair episodes on (scenario, instance, model) and require run status \texttt{COMPLETE} on both sides, giving 167 B3 pairs and 168 CORE pairs. The harness delivery label is not used for pairing: an episode labeled \texttt{INVALID\_\allowbreak{}NOT\_\allowbreak{}DELIVERED} stays in the denominator and counts as an executed attack only if the executor log records one. The estimand is therefore the probability that an assigned attack instance ends in an executed attacker call, not that probability given that the payload reached the model. Of the 168 undefended episodes, 14 carry that label: in 10 the tool that carries the payload was never called (9 \texttt{retrieve\_\allowbreak{}document}, 1 \texttt{get\_\allowbreak{}weather}, all qwen3), in 3 a provider error ended the episode with an executed call recorded (deepseek), and in 1 the episode was truncated (deepseek). Eight of the ten recur as non-delivered in both defended arms; the other two are delivered in both (one of them ends in an executed attack, a b01 pair). Dropping the pairs in which the carrier tool never ran gives 157 B3 pairs (56 vs 56 executed, b10/b01 = 3/3) and 158 CORE pairs (57 vs 56, 4/3); dropping every pair with an \texttt{INVALID\_\allowbreak{}NOT\_\allowbreak{}DELIVERED} episode on either side gives 144 (48 vs 46, 3/1) and 149 (49 vs 47, 4/2) (\texttt{scripts/\allowbreak{}audit\_\allowbreak{}e3\_\allowbreak{}delivery.py}). The qualitative result does not change. The E2 observed-consequence rule (§4.2) is a different eligibility rule and was not applied to E3. \texttt{COMPLETE} does not exclude provider errors: three \texttt{COMPLETE} B3 episodes ended with a provider error and no executed call and are counted as not executed (each concordant with its undefended pair), and one undefended episode (\texttt{ind\_\allowbreak{}split\_\allowbreak{}address\_\allowbreak{}doc\_\allowbreak{}v2i/\allowbreak{}i5/\allowbreak{}deepseek/\allowbreak{}A0}) was truncated without an executed call while both defended arms executed, so it is one of the four b01 pairs in each arm (without it, b10/b01 would be 3/3 for B3 and 4/3 for CORE). Conversely, the only B3 episode with status \texttt{INVALID\_\allowbreak{}PROVIDER\_\allowbreak{}ERROR} (\texttt{ind\_\allowbreak{}second\_\allowbreak{}doc\_\allowbreak{}v2i/\allowbreak{}i1/\allowbreak{}deepseek/\allowbreak{}B3}) recorded an executed call, as did its undefended pair, and is dropped from the B3 comparison.

\section{S2 Tables moved from the main paper}

\textbf{Table 2. Positioning (from the papers\textquotesingle{} own descriptions).} Pathade et al.\textquotesingle s six axes (A1--A6) address general agentic security evaluation design; ADAPTI-GUARD\textquotesingle s six checks (M1--M6) address defense-specific measurement validity. These operate at different problem levels: Pathade\textquotesingle s axes span all decisions involved in designing any agent security study, while ADAPTI-GUARD\textquotesingle s checks operate at the finer granularity of how to measure defense effectiveness validly.

\begin{longtable}[]{@{}
  >{\raggedright\arraybackslash}p{(\columnwidth - 18\tabcolsep) * \real{0.1000}}
  >{\raggedright\arraybackslash}p{(\columnwidth - 18\tabcolsep) * \real{0.1000}}
  >{\raggedright\arraybackslash}p{(\columnwidth - 18\tabcolsep) * \real{0.1000}}
  >{\raggedright\arraybackslash}p{(\columnwidth - 18\tabcolsep) * \real{0.1000}}
  >{\raggedright\arraybackslash}p{(\columnwidth - 18\tabcolsep) * \real{0.1000}}
  >{\raggedright\arraybackslash}p{(\columnwidth - 18\tabcolsep) * \real{0.1000}}
  >{\raggedright\arraybackslash}p{(\columnwidth - 18\tabcolsep) * \real{0.1000}}
  >{\raggedright\arraybackslash}p{(\columnwidth - 18\tabcolsep) * \real{0.1000}}
  >{\raggedright\arraybackslash}p{(\columnwidth - 18\tabcolsep) * \real{0.1000}}
  >{\raggedright\arraybackslash}p{(\columnwidth - 18\tabcolsep) * \real{0.1000}}@{}}
\toprule\noalign{}
\begin{minipage}[b]{\linewidth}\raggedright
\end{minipage} & \begin{minipage}[b]{\linewidth}\raggedright
AgentDojo
\end{minipage} & \begin{minipage}[b]{\linewidth}\raggedright
InjecAgent
\end{minipage} & \begin{minipage}[b]{\linewidth}\raggedright
LLMail-Inject
\end{minipage} & \begin{minipage}[b]{\linewidth}\raggedright
BIPIA
\end{minipage} & \begin{minipage}[b]{\linewidth}\raggedright
Arman et al.
\end{minipage} & \begin{minipage}[b]{\linewidth}\raggedright
Hofer et al.
\end{minipage} & \begin{minipage}[b]{\linewidth}\raggedright
Shaw
\end{minipage} & \begin{minipage}[b]{\linewidth}\raggedright
Pathade et al.
\end{minipage} & \begin{minipage}[b]{\linewidth}\raggedright
this paper
\end{minipage} \\
\midrule\noalign{}
\endhead
\bottomrule\noalign{}
\endlastfoot
success from executed tool or environment state & yes & tool name & yes (tool call) & no (rules and judge) & yes & yes (ground truth) & yes & no (meta-analysis) & yes (executed call) \\
LLM judge in the primary endpoint & no & no & no & partly & no & studied & no & n/a & no (judge compared, not used) \\
human-written attacks & templates & GPT-4-generated, manually revised & yes (839 participants; some LLM-assisted) & authored & reused & automated & n/a & n/a & AI-assistant-authored, partially independent set (§8.3), reported apart from generated items \\
adaptive attackers & manual baselines & follow-up work & yes (participants) & no & no & yes (automated) & n/a & n/a & no \\
evaluation decisions studied as the object & no & no & no & no & one (surface) & one (judge) & several defects & six axes, analytical & six, on the same traces \\
judge versus executed outcome measured & n/a & n/a & n/a & n/a & n/a & yes (precision) & n/a & n/a & yes (counts: 28 of 61 on Track B; 3 and 4 of 61 on Track A) \\
defense applied to the untrusted channel checked & n/a & n/a & n/a & n/a & yes & n/a & n/a & no & yes (§4.7) \\
per-model reporting with validity control and noise floor & partly & partly & partly & partly & per cell & per model & n/a & power analysis only & yes \\
\end{longtable}


\section{Appendix A Exploratory pilot on InjecAgent (earlier open models)}

This pilot preceded the external test of §§5.6 and 6.6. Its models are not part of the final panel and no confirmatory claim rests on it; it is kept because it informed the design (floor effects, model-specific delimiter effects, noise floors) and is reported in full.

To remove the authorship and detector dependence of our own scenarios (M4), we ran cases from InjecAgent (Zhan et al. 2024; external repository, hashes pinned in the run manifest). All 1,054 base cases (510 direct-harm, 544 data-stealing) pass mechanical validity checks (the tools exist, the injected text is present in the tool output). Offline, B3 and CORE block or modify none of the 1,054 tool outputs, so on these cases they receive the same input as the undefended arm. Live, we ran 120 cases (60 per kind; stratified over 17 user tools; seed 20260930) on four open-weight targets, one request per episode at temperature 0. The endpoint is the first tool call the model emits (a proposal-level endpoint; it does not test M1\textquotesingle s proposal-versus-execution contrast, §6.6): an attacker tool or not, with no judge or detector of ours involved. Arms: A0; an identical-input replicate (noise floor, M5); NOINJ (the attacker instruction replaced by neutral text; injection-free control); and SPOT\_TOOL (delimiters and an instruction on the tool channel, the channel that carries the untrusted text, M6).

\begin{longtable}[]{@{}
  >{\raggedright\arraybackslash}p{(\columnwidth - 12\tabcolsep) * \real{0.1429}}
  >{\raggedright\arraybackslash}p{(\columnwidth - 12\tabcolsep) * \real{0.1429}}
  >{\raggedright\arraybackslash}p{(\columnwidth - 12\tabcolsep) * \real{0.1429}}
  >{\raggedright\arraybackslash}p{(\columnwidth - 12\tabcolsep) * \real{0.1429}}
  >{\raggedright\arraybackslash}p{(\columnwidth - 12\tabcolsep) * \real{0.1429}}
  >{\raggedright\arraybackslash}p{(\columnwidth - 12\tabcolsep) * \real{0.1429}}
  >{\raggedright\arraybackslash}p{(\columnwidth - 12\tabcolsep) * \real{0.1429}}@{}}
\toprule\noalign{}
\begin{minipage}[b]{\linewidth}\raggedright
target
\end{minipage} & \begin{minipage}[b]{\linewidth}\raggedright
A0
\end{minipage} & \begin{minipage}[b]{\linewidth}\raggedright
replicate
\end{minipage} & \begin{minipage}[b]{\linewidth}\raggedright
NOINJ
\end{minipage} & \begin{minipage}[b]{\linewidth}\raggedright
SPOT\_TOOL
\end{minipage} & \begin{minipage}[b]{\linewidth}\raggedright
SPOT\_TOOL minus A0, cluster bootstrap 95\% CI
\end{minipage} & \begin{minipage}[b]{\linewidth}\raggedright
A0-only/SPOT-only pairs (exact p, descriptive)
\end{minipage} \\
\midrule\noalign{}
\endhead
\bottomrule\noalign{}
\endlastfoot
qwen-2.5-7b & 24/120 & 28/119 & 0/120 & 26/120 & +0.017 {[}-0.057, 0.086{]} & 10/12 (0.83) \\
llama-3.1-8b & 50/120 & 51/120 & 1/119 & 33/120 & -0.142 {[}-0.237, -0.056{]} & 26/9 (0.006) \\
llama-3.3-70b & 48/120 & 50/120 & 0/119 & 54/120 & +0.050 {[}-0.018, 0.122{]} & 7/13 (0.26) \\
mistral-small-3.2-24b & 5/120 & 5/120 & 0/120 & 5/120 & 0.000 {[}-0.025, 0.024{]} & 1/1 (1.0; floor) \\
\end{longtable}


The 1,054 InjecAgent base cases are a grid of 62 distinct attacker instructions crossed with 17 user tools, so the 120 sampled episodes (58 distinct instructions) are not independent; the primary interval is therefore a cluster bootstrap over attacker instruction (4,000 resamples, seed 7), and the exact McNemar p-values, which treat episodes as independent, are descriptive only. Three observations. (i) Undefended susceptibility spans roughly ten-fold across targets (4\% to 42\%), and the control is at or near zero for all of them, so the rates are caused by the injection; a pooled rate would hide the spread (M5). (ii) Scale within a family did not reduce susceptibility (llama 8B 0.42, 70B 0.40). (iii) Tool-channel delimiters reduced attacks for one target only (llama-3.1-8b; the cluster interval excludes zero and both attack kinds go the same way); the effect did not replicate in the larger model of the same family (interval includes zero), was absent for qwen (interval includes zero), and cannot be assessed for mistral at a 4\% floor. The effect is model-specific, not absent and not general. Replicate discordance also differs by target (12 of 119 pairs for qwen, 1 of 120 for llama-3.1-8b, 0 for mistral), so the noise floor must be measured per model.

Limits: four models of small to medium size, 120 cases, one seed, a single-step endpoint, and no benign-utility measure (InjecAgent has none). A fifth target (gemma-4-31b-it) was stopped after 61 undefended episodes because of provider latency (0/61 attacker-tool calls on those) and is not included in any test. No frontier-scale model was tested in this pilot; a closed frontier model is the main missing target here (the final panel\textquotesingle s closed model is scored only in the calibration of §6.7, not in a defense comparison). Exploratory; no confirmatory claim.

\section{Appendix B Reporting checklist for evaluations of runtime defenses against prompt injection}

Intended for authors and reviewers. Each item has a pass criterion that can be checked from the paper and its artifacts. Items correspond to the candidate checks M1 to M6 of Table 3, derived from this case study. This checklist complements Pathade et al.\textquotesingle s general evaluation-design checklist and is intended for defense-specific evaluation design, not as a universal validation framework.

\begin{longtable}[]{@{}
  >{\raggedright\arraybackslash}p{(\columnwidth - 12\tabcolsep) * \real{0.1429}}
  >{\raggedright\arraybackslash}p{(\columnwidth - 12\tabcolsep) * \real{0.1429}}
  >{\raggedright\arraybackslash}p{(\columnwidth - 12\tabcolsep) * \real{0.1429}}
  >{\raggedright\arraybackslash}p{(\columnwidth - 12\tabcolsep) * \real{0.1429}}
  >{\raggedright\arraybackslash}p{(\columnwidth - 12\tabcolsep) * \real{0.1429}}
  >{\raggedright\arraybackslash}p{(\columnwidth - 12\tabcolsep) * \real{0.1429}}
  >{\raggedright\arraybackslash}p{(\columnwidth - 12\tabcolsep) * \real{0.1429}}@{}}
\toprule\noalign{}
\begin{minipage}[b]{\linewidth}\raggedright
target
\end{minipage} & \begin{minipage}[b]{\linewidth}\raggedright
A0
\end{minipage} & \begin{minipage}[b]{\linewidth}\raggedright
replicate
\end{minipage} & \begin{minipage}[b]{\linewidth}\raggedright
NOINJ
\end{minipage} & \begin{minipage}[b]{\linewidth}\raggedright
SPOT\_TOOL
\end{minipage} & \begin{minipage}[b]{\linewidth}\raggedright
SPOT\_TOOL minus A0, cluster bootstrap 95\% CI
\end{minipage} & \begin{minipage}[b]{\linewidth}\raggedright
A0-only/SPOT-only pairs (exact p, descriptive)
\end{minipage} \\
\midrule\noalign{}
\endhead
\bottomrule\noalign{}
\endlastfoot
qwen-2.5-7b & 24/120 & 28/119 & 0/120 & 26/120 & +0.017 {[}-0.057, 0.086{]} & 10/12 (0.83) \\
llama-3.1-8b & 50/120 & 51/120 & 1/119 & 33/120 & -0.142 {[}-0.237, -0.056{]} & 26/9 (0.006) \\
llama-3.3-70b & 48/120 & 50/120 & 0/119 & 54/120 & +0.050 {[}-0.018, 0.122{]} & 7/13 (0.26) \\
mistral-small-3.2-24b & 5/120 & 5/120 & 0/120 & 5/120 & 0.000 {[}-0.025, 0.024{]} & 1/1 (1.0; floor) \\
\end{longtable}


How to cite: refer to "the M1 to M6 checks" and to "Appendix B of this paper" when stating which items a study satisfies; an item that is not satisfied should be listed as a limitation.

\textbf{Self-assessment of this paper against the checklist.} Satisfied, as we read §4 to §6: items 1, 2, 3, 4, 9, 14, 16. Partly satisfied: items 8 and 10 (the inclusion and admission rules were fixed after the rates were seen, §4.4, §4.5), item 7 (the independent set is frozen by hash but authored by an AI assistant that had read the detector patterns, §8.3), item 12 (instance-cluster bootstrap for E3 only, §8.2), item 13 (checked in the E3 re-run only, §4.7), item 17 (offline numbers regenerate; the live harness and the external-test runner are not released). Not satisfied: item 5 (E2 has no injection-free control; the scenario filter used the either-clause rule of §4.3), item 11 (no undefended replicate in E2 or per model in E3; replicates exist for the external test only), item 15 (no adaptive attacker), item 18 (no external pre-registration). Table 3 gives the status of the checks themselves.

\section{Appendix C Artifacts and data availability}

Unless a row says otherwise the artifact is in the repository. Pack hashes are recorded in the freeze records and \texttt{hashes.sha256} files; the multi-turn run manifests record commit SHAs, not file hashes. Items that are not in the public tree are marked.

\begin{longtable}[]{@{}
  >{\raggedright\arraybackslash}p{(\columnwidth - 6\tabcolsep) * \real{0.2500}}
  >{\raggedright\arraybackslash}p{(\columnwidth - 6\tabcolsep) * \real{0.2500}}
  >{\raggedright\arraybackslash}p{(\columnwidth - 6\tabcolsep) * \real{0.2500}}
  >{\raggedright\arraybackslash}p{(\columnwidth - 6\tabcolsep) * \real{0.2500}}@{}}
\toprule\noalign{}
\begin{minipage}[b]{\linewidth}\raggedright
\#
\end{minipage} & \begin{minipage}[b]{\linewidth}\raggedright
Item
\end{minipage} & \begin{minipage}[b]{\linewidth}\raggedright
Pass criterion
\end{minipage} & \begin{minipage}[b]{\linewidth}\raggedright
Check
\end{minipage} \\
\midrule\noalign{}
\endhead
\bottomrule\noalign{}
\endlastfoot
1 & Endpoint (M1) & success is the executed tool call or the final environment state with argument-level predicates; the proposal rate is reported separately & M1 \\
2 & No judge in the primary endpoint (M1) & if a judge is used, it is compared with the executed outcome and the agreement (kappa, precision, recall) is reported & M1 \\
3 & Blocked payloads (M2) & a block before the target is counted as a defense outcome; results are shown under both labelings & M2 \\
4 & Delivery check & the harness verifies that the payload reached the model and reports how many episodes failed delivery for harness reasons & M2 \\
5 & Validity control (M3) & a control with the injection removed or replaced by neutral text; a model is assessable only if the control is at or near zero & M3 \\
6 & Attacker-controlled effect (M3) & every attack scenario contains an effect the user did not request and the attacker controls & M3 \\
7 & Authorship (M4) & the attack set was written independently of detector or defense development, and frozen (hash) before defenses were run on it; authorship (human, model, mixed) is stated per item & M4 \\
8 & Selection (M4) & any filter used to select attacks (for example "triggered the tool call in an earlier challenge") is stated, and does not use the targets\textquotesingle{} results & M4 \\
9 & Per-model reporting (M5) & results are given per model; pooled rates only as a secondary summary & M5 \\
10 & Floor rule (M5) & a stated rule for models whose undefended rate is too low to assess a defense (for example below 5\% at n \textgreater= 40) & M5 \\
11 & Noise floor (M5) & an undefended replicate shows run-to-run discordance per model & M5 \\
12 & Non-independence (M5) & the unit of analysis accounts for repeated instructions or near-duplicates (cluster bootstrap or equivalent) & M5 \\
13 & Defense channel (M6) & the defense is applied to the channel through which the attack arrives, with a check that it sees the injected text & M6 \\
14 & Utility & benign utility is measured with the same endpoint, with and without the defense & M1 \\
15 & Adaptive attackers & the paper states whether an adaptive attacker was run; if not, no claim of robustness against one & scope \\
16 & Errors & provider failures are reported as failures, never scored as safe; spend and caps are reported & reporting \\
17 & Reproducibility & one command regenerates the offline numbers and figures from committed traces (E1 to E4; external-test numbers are read from committed records); pack hashes are in \texttt{hashes.sha256} files and code commit SHAs are in the run manifests & reporting \\
18 & Pre-registration & hypotheses, endpoint, sample, analysis and floor rule are registered before the confirmatory run; deviations are logged & reporting \\
\end{longtable}


Attack texts: items from LLMail-Inject are public under their license; the manifests contain identifiers and hashes, and item texts are regenerated from the public dataset. Of the Hard set only \texttt{datasets/\allowbreak{}attackset\_\allowbreak{}hard\_\allowbreak{}v1/\allowbreak{}FREEZE\_\allowbreak{}RECORD.json} is in the tree; the manifests it hashes (\texttt{MANIFEST.json}, \texttt{XCHANNEL\_\allowbreak{}MANIFEST\_\allowbreak{}full.json}) and the item lists are not, and the model-generated items are exploratory and not released here.

\section{Appendix D E3 detailed results: per-family and per-model breakdown}

Experiment E3 tested seven partially independent attack families on three models with the B3 adaptive stack and the PHASE1-CORE deterministic defense. Both defenses produced zero blocks across all 336 episodes. PHASE1-CORE applied no transformation (action A0 on every message; the user turn is identical to the undefended run\textquotesingle s in 168 of 168 episodes). B3 applied action A1 on every message and changed the user turn only by trailing punctuation or whitespace in 168 of 168 episodes (\texttt{scripts/\allowbreak{}audit\_\allowbreak{}e3\_\allowbreak{}defense\_\allowbreak{}activity.py}, \texttt{docs/\allowbreak{}research/\allowbreak{}artifacts/\allowbreak{}e3\_\allowbreak{}defense\_\allowbreak{}activity\_\allowbreak{}20261007.json}). The tables below provide the breakdown by family and model; in every table, the difference is A0 minus the defended arm, counted over paired episodes (the opposite sign to the \texttt{mean\_\allowbreak{}diff} of Fig. 3, which is defended minus A0).

Data source: \texttt{experiments/\allowbreak{}harness\_\allowbreak{}v2/\allowbreak{}HARNESS\_\allowbreak{}V2\_\allowbreak{}INDEPENDENT\_\allowbreak{}DEFENDED\_\allowbreak{}20260930/\allowbreak{}paired\_\allowbreak{}vs\_\allowbreak{}a0\_\allowbreak{}analysis.json}

\subsection{E3.1 B3 adaptive stack}

\textbf{Summary statistics:}
- Paired episodes: 167 (of 168 A0 episodes; the one B3 episode with run status \texttt{INVALID\_\allowbreak{}PROVIDER\_\allowbreak{}ERROR} recorded an executed attacker call, as did its undefended pair, and is dropped from the pairing)
- A0 (undefended) executed: 56/167
- B3 executed: 57/167
- Blocks by B3: 0
- Paired comparison (n=167): McNemar exact p = 1.0

\subsection{E3.2 PHASE1-CORE deterministic defense}

\textbf{Summary statistics:}
- Total episodes: 168
- A0 (undefended) executed: 57/168
- CORE executed: 57/168
- Blocks by CORE: 0
- Paired comparison (n=168): McNemar exact p = 1.0

\subsubsection{Per-family and per-model results, B3 and PHASE1-CORE}

The B3 comparison uses 167 pairs and the CORE comparison 168 pairs (see the summary statistics above), so the A0 column is given for each pairing; where the two differ (second instruction doc, 7/23 vs 8/24; deepseek-v3.2, 40/55 vs 41/56) the denominators differ by one.

\begin{longtable}[]{@{}
  >{\raggedright\arraybackslash}p{(\columnwidth - 8\tabcolsep) * \real{0.2000}}
  >{\raggedright\arraybackslash}p{(\columnwidth - 8\tabcolsep) * \real{0.2000}}
  >{\raggedright\arraybackslash}p{(\columnwidth - 8\tabcolsep) * \real{0.2000}}
  >{\raggedright\arraybackslash}p{(\columnwidth - 8\tabcolsep) * \real{0.2000}}
  >{\raggedright\arraybackslash}p{(\columnwidth - 8\tabcolsep) * \real{0.2000}}@{}}
\toprule\noalign{}
\begin{minipage}[b]{\linewidth}\raggedright
artifact
\end{minipage} & \begin{minipage}[b]{\linewidth}\raggedright
content
\end{minipage} & \begin{minipage}[b]{\linewidth}\raggedright
size
\end{minipage} & \begin{minipage}[b]{\linewidth}\raggedright
license or origin
\end{minipage} & \begin{minipage}[b]{\linewidth}\raggedright
hash or pin
\end{minipage} \\
\midrule\noalign{}
\endhead
\bottomrule\noalign{}
\endlastfoot
frozen packs (Tracks A and B) & judge-scored confirmatory packs & 61 + 61 attacks, benign twins & project & SHA-256 in \texttt{datasets/\allowbreak{}frozen/\allowbreak{}*/\allowbreak{}hashes.sha256} \\
harness v2 traces & multi-turn tool episodes with the executor log (\texttt{episodes.jsonl}, manifests, cost summaries); the per-episode trajectories and HTTP streams are not in the tree & 468 (E2), 168 + 336 (E3) episodes & project & commit SHAs in per-run manifests \\
independent scenario set & seven families authored separately from the detector (partially independent; see §8.3); the template file is not in the tree of any branch and is retrievable only by commit SHA (\texttt{e5135a6}) & 7 families × 24 authored instances; 8 per family run (56 instances, 168 episodes with three models) & project & SHA-256 of the template file (§4.4; \texttt{REPRODUCIBILITY.md}) \\
InjecAgent (external) & 62 attacker instructions x 17 contexts & 1,054 base cases & external repository & commit and file hashes in \texttt{datasets/\allowbreak{}external\_\allowbreak{}samples/\allowbreak{}injecagent\_\allowbreak{}phase2\_\allowbreak{}sample.json} \\
Hard set, human-written & tool-call-verified items from LLMail-Inject & 238 items (68 calibration, 170 test) & MIT (Hugging Face \texttt{microsoft/\allowbreak{}llmail-inject-challenge}) & \texttt{datasets/\allowbreak{}attackset\_\allowbreak{}hard\_\allowbreak{}v1/\allowbreak{}FREEZE\_\allowbreak{}RECORD.json} \\
Hard set, model-generated & cross-channel items from a non-target generator & 69 items, exploratory unless validated & project & manifest hash in the freeze record; meta-prompt SHA-256 \\
AgentDojo pairs & fixed seeded pairs for health checks and pilot & 60 pairs of 629 valid & external benchmark & \texttt{datasets/\allowbreak{}external\_\allowbreak{}samples/\allowbreak{}agentdojo\_\allowbreak{}v1\_\allowbreak{}phase2\_\allowbreak{}pairs.json} \\
model panel & registry with pinned prices and settings & 4 open, 1 closed, 2 judges & OpenRouter list & \texttt{configs/\allowbreak{}models\_\allowbreak{}panel\_\allowbreak{}external\_\allowbreak{}v2.yaml} \\
code & public: analysis, figure, ledger and reproduction scripts. Not in the tree of any public branch or tag (present only in unreachable commits fetchable by SHA, \texttt{REPRODUCIBILITY.md}): the live harness (\texttt{src/\allowbreak{}adapti\_\allowbreak{}guard/\allowbreak{}evaluation/\allowbreak{}harness\_\allowbreak{}v2/\allowbreak{}}), the E2/E3 run scripts, the calibration and InjecAgent run scripts, the external-test analysis scripts and the harness test & & project & commit in each manifest (runner commits are retrievable by SHA, not reachable from any branch or tag) \\
reproduction & \texttt{scripts/\allowbreak{}reproduce\_\allowbreak{}negative\_\allowbreak{}result.sh} regenerates the E1 to E4 numbers, ledger and figures offline; the §6.6 and §6.7 numbers (including 23/186 and the intervals of §6.6) are not in \texttt{NUMBERS\_\allowbreak{}LEDGER.md} and are read from committed records (\texttt{experiments/\allowbreak{}external/\allowbreak{}injecagent\_\allowbreak{}registered\_\allowbreak{}20261001/\allowbreak{}ANALYSIS.json}; \texttt{experiments/\allowbreak{}external/\allowbreak{}injecagent\_\allowbreak{}panel\_\allowbreak{}calib\_\allowbreak{}20261001/\allowbreak{}}) and are not regenerated & & & \texttt{NUMBERS\_\allowbreak{}LEDGER.md} with a consistency test \\
\end{longtable}


\begin{longtable}[]{@{}llll@{}}
\toprule\noalign{}
Attack Family & A0 Executed & B3 Executed & Difference (A0 − B3) \\
\midrule\noalign{}
\endhead
\bottomrule\noalign{}
\endlastfoot
authority claim doc & 23/24 & 22/24 & +1 \\
important instructions doc & 9/24 & 8/24 & +1 \\
important instructions weather & 8/24 & 8/24 & 0 \\
second instruction doc & 7/23 & 7/23 & 0 \\
split address doc & 4/24 & 5/24 & −1 \\
unrequested write doc & 5/24 & 7/24 & −2 \\
workflow footer doc & 0/24 & 0/24 & 0 \\
\end{longtable}


\subsection{E3.3 Non-delivered episodes and delivery-restricted pairs}

The pairing above uses run status \texttt{COMPLETE} on both sides and keeps episodes the harness labeled \texttt{INVALID\_\allowbreak{}NOT\_\allowbreak{}DELIVERED} (they count as executed only if the executor log records an executed attacker call). Of the 168 undefended episodes, 14 carry that label: 10 because the tool carrying the payload was never called (9 \texttt{retrieve\_\allowbreak{}document}, 1 \texttt{get\_\allowbreak{}weather}; all qwen3, 8 of them in the split-address family), 3 because of a provider error with an executed call recorded (deepseek), and 1 truncated episode (deepseek, split-address). Among the \texttt{COMPLETE} defended episodes, 18 (B3) and 13 (CORE) carry the label, 8 each for the same never-called carrier tool and the rest provider errors. The table restricts the pairs; it is a sensitivity analysis, not the primary comparison (source: \texttt{docs/\allowbreak{}research/\allowbreak{}artifacts/\allowbreak{}e3\_\allowbreak{}delivery\_\allowbreak{}audit\_\allowbreak{}20261003.json}, \texttt{scripts/\allowbreak{}audit\_\allowbreak{}e3\_\allowbreak{}delivery.py}).

\begin{longtable}[]{@{}llll@{}}
\toprule\noalign{}
Model & A0 Executed & B3 Executed & Difference (A0 − B3) \\
\midrule\noalign{}
\endhead
\bottomrule\noalign{}
\endlastfoot
deepseek-v3.2 & 40/55 & 43/55 & −3 \\
gemma-4-31b-it & 8/56 & 6/56 & +2 \\
qwen3-30b-a3b & 8/56 & 8/56 & 0 \\
\end{longtable}


No restriction produces a significant difference (exact McNemar p ≥ 0.625 in every row); the sample is small and these are descriptive.

\subsection{Interpretation}

Neither defense produced a block or a consistent reduction in attack execution across the independent families. The per-family and per-model differences are small and within the range of run-to-run noise documented in §6.3 (8 of 168 paired outcomes differed at temperature 0 for CORE; 7 of 167 for B3), so differences of this size cannot be attributed to a defense.

\end{document}
